% from ega4-18-fr.tex
\oldpage[IV-4]{139}

est le radical de ${}^h\!A$ et que l'application canonique $A/\mathfrak r\to{}^h\!A/\mathfrak r.{}^h\!A$ est bijective. Comme le séparé complété $\widehat A$ de $A$ pour la topologie $\mathfrak r$-préadique est le produit des séparés complétés $\widehat A_{\mathfrak m_i}$ (Bourbaki, \emph{Alg. comm.}, chap. III, \textsection~2, n\textsuperscript{o}~13, prop. 18) l'homomorphisme canonique $\widehat A\to({}^h\!A)\widehat{\phantom{x}}$ est bijectif par (18.6.6, (iv)), et il est clair par (18.6.6, (v)) que pour que ${}^h\!A$ soit noethérien, il faut et il suffit que $A$ le soit.

Pour obtenir l'analogue de la propriété universelle (18.6.6, (ii)), convenons, lorsque $A$ et $B$ sont deux anneaux semi-locaux de radicaux respectifs $\mathfrak r$, $\mathfrak s$, d'appeler \emph{homomorphisme semi-local} de $A$ dans $B$ tout homomorphisme $\varphi$ tel que $\varphi(\mathfrak r)\subset\mathfrak s$. Dans le cas particulier où $B$ est un \emph{produit} d'anneaux locaux $B_j$ ($1\leq j\leq n$), la donnée d'un tel homomorphisme revient à celle de ses projections $\varphi_j:A\to B_j$, qui sont des homomorphismes soumis à la seule condition que $\varphi_j(\mathfrak r)\subset\mathfrak n_j$, où $\mathfrak n_j$ est l'idéal maximal de $B_j$. D'ailleurs, si $\mathfrak m_i$ ($1\leq i\leq m$) sont les idéaux maximaux de $A$, l'idéal premier $\varphi_j^{-1}(\mathfrak n_j)$ doit contenir l'un des $\mathfrak m_i$ puisqu'il contient leur intersection $\mathfrak r$, donc il est égal à l'un des $\mathfrak m_i$ et $\varphi_j$ se factorise en $A\to A_{\mathfrak m_i}\xrightarrow{\psi_{ij}}B_j$, où $\psi_{ij}$ est un homomorphisme \emph{local}. Si $\Hom.\operatorname{sloc}(A,B)$ désigne l'ensemble des homomorphismes semi-locaux de $A$ dans $B$, on peut donc identifier cet ensemble à
\[
  \prod_j\left(\bigcup_i\Hom.\operatorname{loc}(A_{\mathfrak m_i},B_j)\right)
\]
dans le cas considéré; comme un anneau semi-local hensélien est produit direct d'anneaux locaux henséliens, on voit que la propriété universelle (18.6.6, (ii)) est encore valable pour un anneau semi-local $A$, lorsqu'on y remplace les homomorphismes locaux par les homomorphismes semi-locaux.
\end{env}

\begin{proposition}[18.6.8]
\label{IV.18.6.8}
Soient $A$ un anneau semi-local, $B$ une $A$-algèbre semi-locale et entière sur $A$. Alors $B\otimes_A({}^h\!A)$ est une $B$-algèbre semi-locale isomorphe à ${}^h\!B$.
\end{proposition}

Comme $B\otimes_A({}^h\!A)$ est entier au-dessus de ${}^h\!A$, chacun de ses idéaux maximaux est au-dessus d'un des idéaux maximaux de ${}^h\!A$, donc au-dessus d'un idéal maximal de $A$, et comme $B$ est entière sur $A$, les idéaux premiers de $B$ au-dessus d'un idéal maximal de $A$ sont des idéaux maximaux de $B$, si bien que finalement un idéal maximal de $B\otimes_A({}^h\!A)$ a pour projections dans $\Spec({}^h\!A)$ et $\Spec(B)$ des idéaux maximaux; compte tenu de (18.6.6, (iii)) et de (I, 3.4.9), on conclut que l'anneau $B\otimes_A({}^h\!A)$ est semi-local. En outre, pour tout anneau local hensélien $C$, $\Hom.\operatorname{sloc}(B\otimes_A({}^h\!A),C)$ est en correspondance biunivoque avec l'ensemble des couples $(\varphi,\psi)$ d'homomorphismes semi-locaux $\varphi:B\to C$, $\psi:{}^h\!A\to C$ tels que les composés
\[
  A\longrightarrow B\xrightarrow{\varphi}C,\qquad
  A\longrightarrow{}^h\!A\xrightarrow{\psi}C
\]
soient égaux. Mais en vertu de la correspondance biunivoque entre $\Hom.\operatorname{sloc}(A,C)$ et $\Hom.\operatorname{sloc}({}^h\!A,C)$ on voit que pour tout $\varphi\in\Hom.\operatorname{sloc}(B,C)$ il existe un $\psi$ et un seul ayant la propriété précédente, donc l'application
\[
  \Hom.\operatorname{sloc}(B\otimes_A({}^h\!A),C)\longrightarrow\Hom.\operatorname{sloc}(B,C)
\]
est bijective, ce qui prouve la proposition en vertu de l'unicité de la solution d'un problème universel.

\begin{theorem}[18.6.9]
\label{IV.18.6.9}
Soit $A$ un anneau semi-local.
\begin{enumerate}
  \item[(i)] Pour que ${}^h\!A$ soit réduit (resp. normal), il faut et il suffit que $A$ le soit.
  \item[(ii)] Supposons $A$ noethérien. Alors, pour tout idéal premier $\mathfrak p$ de $A$, l'anneau $({}^h\!A)_{\mathfrak p}/\mathfrak p.({}^h\!A)_{\mathfrak p}$ est composé direct d'un nombre fini d'extensions algébriques séparables de $k(\mathfrak p)$ (ce qui entraîne que les fibres du morphisme canonique $\Spec({}^h\!A)\to\Spec(A)$ sont \emph{géométriquement régulières}).
\end{enumerate}
\end{theorem}

